% !TEX root = ../main.tex
% ---------------------------------------------------------------
% GENERATED FILE - DO NOT EDIT.
% Produced by docs/tools/render_reference.py from the repository source.
% Regenerate with: make docs
% ---------------------------------------------------------------

\apibreak
\section{\texttt{QueryRequest}}
\label{class:feedback_intelligence_agent.schemas.QueryRequest}\label{schema:QueryRequest}\index{Classes!QueryRequest}

API request for the `/query` endpoint.
\apifield{Inherits}\texttt{BaseModel}
\apifield{Fields}
\begin{arglist}
  \item[\texttt{question}] \texttt{str}, required, \texttt{min\_\allowbreak{}length=3}
  \item[\texttt{top\_\allowbreak{}k}] \texttt{int}, default \texttt{4}, \texttt{ge=1,\allowbreak{} le=12}
  \item[\texttt{tenant\_\allowbreak{}id}] \texttt{str \textbar{} None}, default \texttt{None}, \texttt{min\_\allowbreak{}length=1}
  \item[\texttt{customer\_\allowbreak{}segment}] \texttt{str \textbar{} None}, default \texttt{None}, \texttt{min\_\allowbreak{}length=1}
  \item[\texttt{channel}] \texttt{FeedbackChannel \textbar{} None}, default \texttt{None}
  \item[\texttt{min\_\allowbreak{}rating}] \texttt{int \textbar{} None}, default \texttt{None}, \texttt{ge=1,\allowbreak{} le=5}
  \item[\texttt{max\_\allowbreak{}rating}] \texttt{int \textbar{} None}, default \texttt{None}, \texttt{ge=1,\allowbreak{} le=5}
  \item[\texttt{created\_\allowbreak{}after}] \texttt{datetime \textbar{} None}, default \texttt{None}
  \item[\texttt{created\_\allowbreak{}before}] \texttt{datetime \textbar{} None}, default \texttt{None}
\end{arglist}
\apifield{Source}\texttt{src/\allowbreak{}feedback\_\allowbreak{}intelligence\_\allowbreak{}agent/\allowbreak{}schemas.\allowbreak{}py:170}~(\href{https://github.com/DiogoRibeiro7/feedback-intelligence-agent/blob/936886da072fa7025a57c5517b9c90772bedae09/src/feedback_intelligence_agent/schemas.py#L170}{online}), pydantic model, module \cref{module:feedback_intelligence_agent.schemas}

\subsection{\texttt{validate\_filter\_ranges}}
\label{method:feedback_intelligence_agent.schemas.QueryRequest.validate_filter_ranges}\index{Methods!validate\_filter\_ranges}
Validate metadata filters at request parsing time.
\apifield{Usage}
\begin{lstlisting}[style=usage, language=Python]
validate_filter_ranges(self) -> QueryRequest
\end{lstlisting}
\apifield{Value}\texttt{QueryRequest}
\apifield{Source}\texttt{src/\allowbreak{}feedback\_\allowbreak{}intelligence\_\allowbreak{}agent/\allowbreak{}schemas.\allowbreak{}py:184}~(\href{https://github.com/DiogoRibeiro7/feedback-intelligence-agent/blob/936886da072fa7025a57c5517b9c90772bedae09/src/feedback_intelligence_agent/schemas.py#L184}{online})

\subsection{\texttt{metadata\_filters}}
\label{method:feedback_intelligence_agent.schemas.QueryRequest.metadata_filters}\index{Methods!metadata\_filters}
Return the request's metadata filters, or None when no filter is set.
\apifield{Usage}
\begin{lstlisting}[style=usage, language=Python]
metadata_filters(self) -> MetadataFilters | None
\end{lstlisting}
\apifield{Value}\texttt{MetadataFilters \textbar{} None}
\apifield{Source}\texttt{src/\allowbreak{}feedback\_\allowbreak{}intelligence\_\allowbreak{}agent/\allowbreak{}schemas.\allowbreak{}py:189}~(\href{https://github.com/DiogoRibeiro7/feedback-intelligence-agent/blob/936886da072fa7025a57c5517b9c90772bedae09/src/feedback_intelligence_agent/schemas.py#L189}{online})
